\chapter{Lesson 6: Hobbies}

\section{Vocabulary}

\begin{center}
\begin{tabular}{llll}
\toprule
Characters & Pinyin & Part of speech & Meaning \\
\midrule
\zh{爱好} & àihào & noun & hobby \\
\zh{周末} & zhōumò & noun & weekend \\
\zh{打球} & dǎ qiú & verb-object & to play ball \\
\zh{看电视} & kàn diànshì & verb-object & to watch TV \\
\zh{看电影} & kàn diànyǐng & verb-object & to see a movie \\
\zh{唱歌} & chàng gē & verb-object & to sing (songs) \\
\zh{跳舞} & tiào wǔ & verb-object & to dance \\
\zh{听音乐} & tīng yīnyuè & verb-object & to listen to music \\
\zh{看书} & kàn shū & verb-object & to read books \\
\zh{游泳} & yóu yǒng & verb-object & to swim \\
\zh{玩游戏} & wán yóuxì & verb-object & to play (video) games \\
\zh{对} & duì & adjective & right, correct \\
\zh{有的时候} & yǒude shíhou & phrase & sometimes \\
\zh{常常} & chángcháng & adverb & often \\
\zh{那} & nà & conjunction & in that case, then \\
\zh{去} & qù & verb & to go \\
\zh{觉得} & juéde & verb & to feel, to think \\
\zh{有意思} & yǒu yìsi & adjective & interesting \\
\zh{没意思} & méi yìsi & adjective & boring \\
\zh{只} & zhǐ & adverb & only \\
\zh{睡觉} & shuì jiào & verb-object & to sleep \\
\zh{算了} & suàn le & phrase & forget it, never mind \\
\zh{找} & zhǎo & verb & to look for \\
\zh{别人} & biérén & pronoun & others, other people \\
\zh{好久} & hǎo jiǔ & phrase & a long time \\
\bottomrule
\end{tabular}
\end{center}

\section{Verb--object compounds}

Many activities that are a single verb in English are \emph{verb + object} in Chinese:
\zh{唱歌} is literally ``sing song'', \zh{跳舞} ``jump dance'', \zh{睡觉} ``sleep a
sleep'', \zh{吃饭} ``eat rice''.

\begin{important}[Consequences of verb--object structure]
\begin{itemize}
  \item The object can be replaced by a more specific one: \zh{唱中文歌} ``sing Chinese
        songs'', \zh{看中国电影}.
  \item Things can be inserted between the two parts (measure phrases, \zh{了}, \zh{过},
        complements): \zh{睡了一个觉} ``had a sleep''. We will see much more of this in
        CHINA 102R.
  \item A verb--object compound cannot take a second object: ``to dance with her'' is
        \zh{跟她跳舞}, not *\zh{跳舞她}.
\end{itemize}
\end{important}

\section{Grammar}

\subsection{\zh{喜欢} + verb phrase}

\zh{喜欢} can be followed by a noun (\zh{我喜欢音乐}) or by a whole verb phrase.
\begin{itemize}
  \item \zh{我喜欢看电影。} I like watching movies.
  \item \zh{你喜欢不喜欢跳舞？} Do you like dancing?
  \item \zh{我不太喜欢唱歌。} I don't really like singing. (\zh{不太} ``not very'')
\end{itemize}

\subsection{Affirmative--negative (A-not-A) questions}

Another way to ask a yes/no question is to say the verb or adjective twice, positive
then negative.

\begin{center}
\begin{tabular}{lll}
\toprule
With \zh{吗} & A-not-A & Meaning \\
\midrule
\zh{你忙吗？} & \zh{你忙不忙？} & Are you busy? \\
\zh{他是老师吗？} & \zh{他是不是老师？} & Is he a teacher? \\
\zh{你有哥哥吗？} & \zh{你有没有哥哥？} & Do you have an older brother? \\
\zh{你喜欢游泳吗？} & \zh{你喜欢不喜欢游泳？} / \zh{你喜不喜欢游泳？} & Do you like swimming? \\
\bottomrule
\end{tabular}
\end{center}

\begin{center}
\begin{tikzpicture}[node distance=0.6cm and 0.9cm, b/.style={draw, rounded corners, align=center, font=\small, fill=blue!6, text width=3.6cm}, d/.style={draw, diamond, aspect=2.4, align=center, font=\small, fill=yellow!20, inner sep=1pt}]
  \node[d] (q) {yes/no question?};
  \node[d, below=of q] (adv) {adverb like\\ \zh{也}, \zh{都}, \zh{很}, \zh{常常}\\ before the verb?};
  \node[b, right=1.2cm of adv] (ma) {use \zh{吗} only:\\ \zh{你也喜欢跳舞吗？}};
  \node[b, below=of adv] (both) {either form works:\\ \zh{你忙吗？} / \zh{你忙不忙？}};
  \draw[->, thick] (q) -- node[right, font=\scriptsize] {yes} (adv);
  \draw[->, thick] (adv) -- node[above, font=\scriptsize] {yes} (ma);
  \draw[->, thick] (adv) -- node[right, font=\scriptsize] {no} (both);
\end{tikzpicture}
\end{center}

\begin{important}[Restrictions on A-not-A]
\begin{itemize}
  \item Never add \zh{吗} to an A-not-A question: *\zh{你忙不忙吗？}
  \item Adverbs such as \zh{也, 都, 很, 常常} cannot be used with A-not-A:
        *\zh{你也忙不忙？} Use \zh{吗}: \zh{你也忙吗？}
  \item For \zh{有} the negative is \zh{没}: \zh{有没有}.
  \item With two-syllable verbs, the short form AB-not-AB $\rightarrow$ A-not-AB is
        common in speech: \zh{喜不喜欢}, \zh{认不认识}.
\end{itemize}
\end{important}

\subsection{Going somewhere to do something: \zh{去} + place + verb}

\zh{去} + place + verb phrase expresses purpose: ``go to X to do Y''.
\begin{itemize}
  \item \zh{我们去电影院看电影吧。} Let's go to the cinema to see a movie.
  \item \zh{他常常去图书馆看书。} He often goes to the library to read.
  \item \zh{周末你想去哪儿玩儿？} Where do you want to go for fun this weekend?
\end{itemize}

\subsection{\zh{那(么)} ``in that case''}

\zh{那} at the start of a reply picks up what the other person said and draws a
conclusion.
\begin{itemize}
  \item \textbf{A:} \zh{我不喜欢看电影。}
  \item \textbf{B:} \zh{那我们去跳舞，怎么样？} Then let's go dancing, how about it?
\end{itemize}

\subsection{\zh{觉得}}

\zh{觉得} ``to feel / think'' gives an opinion, often about whether something is
\zh{有意思}.
\begin{itemize}
  \item \zh{我觉得中文很有意思。} I think Chinese is interesting.
  \item \zh{你觉得这个电影怎么样？}---\zh{我觉得没有意思。}
\end{itemize}

\subsection{Frequency adverbs}

\zh{常常} ``often'', \zh{有的时候} ``sometimes'', \zh{不常} ``not often'' all go before
the verb (\zh{有的时候} can also start the sentence).
\begin{itemize}
  \item \zh{我周末常常打球，有的时候也看书。}
  \item \zh{他不常唱歌。}
\end{itemize}

\begin{aside}
\zh{好久不见！} \emph{Hǎojiǔ bú jiàn!} ``Long time no see!'' is a genuine Chinese
expression, and English probably borrowed it in the 19th century.
Weekend culture: KTV (karaoke rooms) is a hugely popular way for Chinese students to
spend a Friday night; the Chinese Students and Scholars Association at Waterloo held
KTV nights and mid-autumn festival parties that the class was encouraged to attend.
\end{aside}

\section{Model conversation}

\begin{itemize}
  \item \textbf{A:} \zh{小高，你周末喜欢做什么？}
  \item \textbf{B:} \zh{我喜欢打球、听音乐。你呢？}
  \item \textbf{A:} \zh{我喜欢看电影。这个周末我们去看电影，怎么样？}
  \item \textbf{B:} \zh{好啊。看中国电影还是美国电影？}
  \item \textbf{A:} \zh{看中国电影吧。我觉得中国电影很有意思。}
  \item \textbf{B:} \zh{好。看电影以后我请你吃饭。}
  \item \textbf{A:} \zh{太好了！}
\end{itemize}

\section*{Review}
\begin{itemize}
  \item Many activities are verb--object compounds (\zh{唱歌, 跳舞, 睡觉}); the object
        can be specified and other words can be inserted.
  \item \zh{喜欢} + verb phrase.
  \item A-not-A questions: no \zh{吗}, no adverbs like \zh{也/都/很}; \zh{有没有}.
  \item \zh{去} + place + verb = go somewhere to do something.
  \item \zh{那} ``in that case''; \zh{觉得} for opinions.
  \item Frequency adverbs go before the verb.
\end{itemize}
